\section{资源与成本透明}

Yuandong Tian 多次提醒：统一框架的核心不是继续单点堆大模型，而是把有限资源分配到搜索、对话采样、评估与优化这些互相耦合的环节中。

\subsection{训练 / 推理成本对比}

\begin{table}[H]
\centering
\begin{tabular}{lrrr}
\toprule
组件 & 主要消耗 & 典型参数量 & 观测指标 \\
\midrule
Search Former & GPU hours × trace length & 10-20M & Trace fidelity、diversity coverage \\
APAC agent & Token cost × 4-6 rounds & 500M-1B & Dialogue throughput、judge reward \\
DSPy solver & Solver iterations × latents & 5-10M & Constraint slip rate、latents convergence \\
\bottomrule
\end{tabular}
\caption{Compound AI 各模块资源分布}
\end{table}

\begin{importantbox}{分配策略}
1）在 search trace 上节省 token，但保留关键 cost info；2）在 APAC 上用 lightweight fine-tuning 快速迭代 judgment；3）在 DSPy 上追踪 solver steps，避免 gradient explosion。这样的资源组合比纯 scaling 更省钱。
\end{importantbox}

\subsection{Inference latency 管控}

Inference 时，trace generation → APAC dialogue → DSPy solver 三个阶段串联，整体 latency 受到最慢环节限制。实际部署采用如下策略：
\begin{itemize}
    \item Search trace precompute + cache heatmap；
    \item APAC 设定最小轮数 3，超过 5 轮则触发 fallback；
    \item DSPy solver use warm start，即 reuse 上一次 \(C\) 初始化。
\end{itemize}

\begin{knowledgebox}{Latency guardrails}
在推理 pipeline 中设置 \texttt{max\_trace\_steps}、\texttt{max\_dialog\_rounds}、\texttt{max\_solver\_iters}，避免某个环节因为异常而拖垮全链路。监控必须包括 queue length 和 rollout duration。
\end{knowledgebox}

\subsection{本章小结}

分解资源（trace / agent / solver）可以更准地衡量 Compound AI 的成本收益，同时在 inference 中设置 guardrail 保障 latency。

